% Options for packages loaded elsewhere
\PassOptionsToPackage{unicode}{hyperref}
\PassOptionsToPackage{hyphens}{url}
\PassOptionsToPackage{space}{xeCJK}
\documentclass[
]{article}
\usepackage{xcolor}
\usepackage{amsmath,amssymb}
\setcounter{secnumdepth}{-\maxdimen} % remove section numbering
\usepackage{iftex}
\ifPDFTeX
  \usepackage[T1]{fontenc}
  \usepackage[utf8]{inputenc}
  \usepackage{textcomp} % provide euro and other symbols
\else % if luatex or xetex
  \usepackage{unicode-math} % this also loads fontspec
  \defaultfontfeatures{Scale=MatchLowercase}
  \defaultfontfeatures[\rmfamily]{Ligatures=TeX,Scale=1}
\fi
\usepackage{lmodern}
\ifPDFTeX\else
  % xetex/luatex font selection
  \setmainfont[]{Times New Roman}
  \setsansfont[]{SimHei}
  \setmonofont[]{SimKai}
  \ifXeTeX
    \usepackage{xeCJK}
    \setCJKmainfont[]{SimSun}
  \fi
  \ifLuaTeX
    \usepackage[]{luatexja-fontspec}
    \setmainjfont[]{SimSun}
  \fi
\fi
% Use upquote if available, for straight quotes in verbatim environments
\IfFileExists{upquote.sty}{\usepackage{upquote}}{}
\IfFileExists{microtype.sty}{% use microtype if available
  \usepackage[]{microtype}
  \UseMicrotypeSet[protrusion]{basicmath} % disable protrusion for tt fonts
}{}
\makeatletter
\@ifundefined{KOMAClassName}{% if non-KOMA class
  \IfFileExists{parskip.sty}{%
    \usepackage{parskip}
  }{% else
    \setlength{\parindent}{0pt}
    \setlength{\parskip}{6pt plus 2pt minus 1pt}}
}{% if KOMA class
  \KOMAoptions{parskip=half}}
\makeatother
\usepackage{longtable,booktabs,array}
\newcounter{none} % for unnumbered tables
\usepackage{calc} % for calculating minipage widths
% Correct order of tables after \paragraph or \subparagraph
\usepackage{etoolbox}
\makeatletter
\patchcmd\longtable{\par}{\if@noskipsec\mbox{}\fi\par}{}{}
\makeatother
% Allow footnotes in longtable head/foot
\IfFileExists{footnotehyper.sty}{\usepackage{footnotehyper}}{\usepackage{footnote}}
\makesavenoteenv{longtable}
\ifLuaTeX
  \usepackage{luacolor}
  \usepackage[soul]{lua-ul}
\else
  \usepackage{soul}
  \ifXeTeX
    % soul's \st doesn't work for CJK:
    \usepackage{xeCJKfntef}
    \renewcommand{\st}[1]{\sout{#1}}
  \fi
\fi
\setlength{\emergencystretch}{3em} % prevent overfull lines
\providecommand{\tightlist}{%
  \setlength{\itemsep}{0pt}\setlength{\parskip}{0pt}}
\usepackage{bookmark}
\IfFileExists{xurl.sty}{\usepackage{xurl}}{} % add URL line breaks if available
\urlstyle{same}
\hypersetup{
  hidelinks,
  pdfcreator={LaTeX via pandoc}}

\author{}
\date{}

\begin{document}

湖北第二师范学院

本科毕业论文（设计）

文献综述

　　　　　

{\def\LTcaptype{none} % do not increment counter
\begin{longtable}[]{@{}
  >{\raggedright\arraybackslash}p{(\linewidth - 2\tabcolsep) * \real{0.3091}}
  >{\raggedright\arraybackslash}p{(\linewidth - 2\tabcolsep) * \real{0.6909}}@{}}
\toprule\noalign{}
\begin{minipage}[b]{\linewidth}\centering
\textbf{论文(设计)题目}:
\end{minipage} & \begin{minipage}[b]{\linewidth}\raggedright
\ul{基于粒子滤波算法的多目标跟踪方法的研究与实现}

\ul{峰粉}
\end{minipage} \\
\midrule\noalign{}
\endhead
\bottomrule\noalign{}
\endlastfoot
\end{longtable}
}

　　　　　

{\def\LTcaptype{none} % do not increment counter
\begin{longtable}[]{@{}
  >{\raggedright\arraybackslash}p{(\linewidth - 2\tabcolsep) * \real{0.2905}}
  >{\centering\arraybackslash}p{(\linewidth - 2\tabcolsep) * \real{0.6665}}@{}}
\toprule\noalign{}
\begin{minipage}[b]{\linewidth}\raggedright
\textbf{姓 名:}
\end{minipage} & \begin{minipage}[b]{\linewidth}\centering
余慧君
\end{minipage} \\
\midrule\noalign{}
\endhead
\bottomrule\noalign{}
\endlastfoot
\textbf{学 号:} & 1750341028 \\
\textbf{年 级:} & 2017 级 \\
\textbf{学 院:} & 计算机学院 \\
\textbf{专 业 名 称:} & 计算机科学与技术 \\
\textbf{指导教师姓名}: & 王海军；李云雪 \\
\textbf{指导教师职称}: & 教授；高级工程师 \\
\end{longtable}
}

填表时间： 年 月 日

\textbf{1. 概述}

正文文字正文文字正文文字正文文字正文文字正文文字正文文字正文文字正文文正文文字正文文字正文文字正文文字正文文字正文文字。

正文文字正文文字正文文字正文文字正文文字正文文字正文文字正文文字正文文正文文字正文文字正文文字正文文字正文文字正文文字。

正文文字正文文字正文文字正文文字正文文字正文文字正文文字正文文字正文文正文文字正文文字正文文字正文文字正文文字正文文字。

\textbf{2. 主题}

正文文字正文文字，正文文字正文文字\textsuperscript{{[}1{]}}，正文文字正文文字\textsuperscript{{[}2{]}}。正文文字正文文字\textsuperscript{{[}3{]}}，正文文字正文文字正文\textsuperscript{{[}4{]}}，正文文字正文文字正文文字\textsuperscript{{[}5{]}}。正文文字正文文字，正文文字正文文字正文文字正文文字正文文字正文文字正文文字正文文字\textsuperscript{{[}6{]}}。正文文字正文文字正文文字正文文字正文文字。

正文文字正文文字\textsuperscript{{[}7{]}}，正文文字正文文字正文文字正文文字\textsuperscript{{[}8{]}}。正文文字正文文字正文文字正文文字正文文字正文文字正文文字\textsuperscript{{[}9{]}}，正文文字正文文字正文文字正文文字正文文字正文文字正文文字\textsuperscript{{[}10{]}}，正文文字正文文字正文文字。

正文文字正文文字正文文字\textsuperscript{{[}11{]}}，正文文字正文文字正文文字正文文字正文文字正文文字\textsuperscript{{[}12{]}}。正文文字正文文字正文文字\textsuperscript{{[}13{]}}，正文文字正文文字正文文字正文文字正文文字正文文字正文文字正文文字正文文字正文文字正文文字正文文字。

正文文字正文文字\textsuperscript{{[}14{]}}，正文文字正文文字正文文字\textsuperscript{{[}15{]}}。正文文字正文文字正文文字正文文字\textsuperscript{{[}16{]}}，正文文字正文文字正文文字\textsuperscript{{[}17{]}}。正文文字正文文字正文文字正文文字正文文字正文文字正文文字\textsuperscript{{[}18{]}}，正文文字正文文字正文。

\textbf{3. 总结}

正文文字正文文字正文文字正文文字正文文字正文文字正文文字正文文字正文文字正文文字正文文字正文文字正文文字正文文字正文文字正文文字正文文字正文文字正文文字正文文字正文文字正文文字正文文字正文文字。

正文文字正文文字正文文字正文文字正文文字正文文字正文文字正文文字正文文字正文文字正文文字正文文字正文文字正文文字正文文字正文文字正文文字正文文字正文文字正文文字正文文字正文文字正文文字正文文字。

正文文字正文文字正文文字正文文字正文文字正文文字正文文字正文文字正文文字正文文字正文文字正文文字正文文字正文文字正文文字正文文字正文文字正文文字正文文字正文文字正文文字正文文字正文文字正文文字。

\textbf{4. 参考文献}

{[}1{]} 杨亚男, 付春玲. 视频目标跟踪算法综述{[}J{]}. 科技资讯, 2018,
16(02): 14-15.

{[}2{]} 徐群, 李家辉, 陈琛, 等.
人工智能技术的智能视频监控系统的相关探究{[}J{]}. 山东工业技术, 2019,
12(09): 131-134.

{[}3{]} 杨艳. 智能交通中运动目标检测与跟踪算法研究{[}D{]}. 昆明:
昆明理工大学, 2016.

{[}4{]} 郜璐璐. 浅析基于多目标跟踪的医学图像分析{[}J{]}. 影像技术. 2014,
29(3): 35-36.

{[}5{]} 张剑华. 均值漂移算法在视频运动目标跟踪中的应用研究{[}D{]}. 西安:
西安电子科技大学, 2015.

{[}6{]} 周梦姣. 基于多特征融合的目标跟踪算法{[}D{]}. 西安:
西安电子科技大学, 2017.

{[}7{]} Han Z, Zhang R, L Wen, et al. Moving Object Tracking Method
Based on Improved Camshift Algorithm{[}J{]}. International Conference on
Industrial Informatics-computing Technology, 2016, pp: 91-95.

{[}8{]} Guo XJ, Wan L, Liu F, et al. An Algorithm for Small Space Target
Tracking Based on Extended Kalman Filter{[}J{]}. Electronics Optics \&
Control, 2016, 4(2): 36-41.

{[}9{]} 权义萍, 金鑫, 张蕾, 等. 基于 Mean-Shift
的卡尔曼粒子滤波车辆跟踪算法{[}J{]}. 南京: 南京工业大学学报, 2014, 3(1):
12-16.

{[}10{]} 宁纪锋. 图像分割和目标跟踪中的若干问题研究{[}D{]}. 西安:
西安电子科技大学, 2009.

{[}11{]} 朱胜利. Mean-Shift 及相关算法在视频跟踪中的研究{[}D{]}. 杭州:
浙江大学, 2006.

{[}12{]} 龙煜. 运动目标检测和遮挡条件下目标跟踪的研究{[}D{]}. 武汉:
武汉科技大学, 2018.

{[}13{]} SHA Musavi, BS Chowdhry, J Bhatti. Object Tracking based on
ActiveContour Modeling{[}C{]}. In: Proceedings of International
Conference on Wireless Communications. 2014, pp: 1-5.

{[}14{]} Weiwu Ren, Xiao Chen, Xiaoming Wang, et al. A Novel Optical
Flow Algorithm Based on Bionic Features for Robust Tracing{[}M{]}.
Germany: Springer International Publishing. 2015, pp: 235-240.

{[}15{]} Dušan Krokavec and Anna Filasová. Fault Residuals Based on
Distributed Discrete-Time Linear Kalman Filtering{[}M{]}. London:
IntechOpen, 2018. pp: 110-125.

{[}16{]} 贺超. 基于图像频域分析显著目标检测算法研究{[}D{]}.
济南:山东大学, 2017.

{[}17{]} 黄金海. 背景帧间差分法的移动目标跟踪研究{[}J{]}. 中国仪器仪表,
2019(01): 62-65.

{[}18{]} 王奎奎, 玉振明. 融合背景减法和帧差法的运动目标检测{[}J{]}.
电视技术, 2015, 39(24): 94-99.

\subsubsection{（要求：15篇以上参考文献，近5年特别是近3年文献数量充足，至少3篇外文参考文献，正文中需要上标引用参考文献。）}\label{ux8981ux6c4215ux7bc7ux4ee5ux4e0aux53c2ux8003ux6587ux732eux8fd15ux5e74ux7279ux522bux662fux8fd13ux5e74ux6587ux732eux6570ux91cfux5145ux8db3ux81f3ux5c113ux7bc7ux5916ux6587ux53c2ux8003ux6587ux732eux6b63ux6587ux4e2dux9700ux8981ux4e0aux6807ux5f15ux7528ux53c2ux8003ux6587ux732e}

\end{document}
